% \iffalse meta-comment
%
% hit-msword.dtx 是把 Word 页面设置换算成 LaTeX 版面参数
%
% \fi
%
% \section{Word 页面设置的换算}
%
% \changes{v3.2a}{2026/08/13}{加 \cs{hit@msword@layout:n}，按 Word 的页面设置算出
%   版心、行距、首行位置与字距}
% 学校发的范例是 Word 排的，规范里写的也是 Word 的话：「页边距上 3.8 厘米」「文档网格
% 每页 33 行、间距 19.55 磅」。这些量与 \hologo{LaTeX} 的 \cs{textheight}、\cs{baselineskip}、
% \cs{topskip} 不是一一对应的，中间隔着 Word 自己的几条规则。以前是把换算结果硬写进
% \cs{geometry} 里，一份变体一组数，改一个值要连带手算三四个。
%
% \cs{hit@msword@layout:n}\marg{键值表} 收 Word 那边的原话，算出 LaTeX 这边的参数：
% \begin{latex}
% \hit@msword@layout:n
%   {
%     margin-top = 3.8cm , margin-bottom = 3cm ,
%     margin-left = 3cm , margin-right = 3cm ,
%     header = 3cm , footer = 2.3cm ,
%     lines-per-page = 33 , line-pitch = 19.55bp ,
%     chars-per-line = 34 , char-pitch = 12.45bp ,
%     size = 12bp , head = 5mm ,
%   }
% \end{latex}
%
% \subsection{Word 的几条规则}
%
% 下面每一条都是从范例与专门造的对照文档里量出来的，不是从文档里读来的。量法是拿
% \texttt{mutool} 读 PDF 内容流里的 \texttt{Tm}，取到的是 Word 自己写下的基线坐标。
%
% 先把「磅」这个字定死。Word 界面上的磅是一英寸的七十二分之一，
% 也就是 \hologo{LaTeX} 的 \texttt{bp}；\hologo{LaTeX} 自己的 \texttt{pt} 是一英寸的
% 72.27 分之一，两者差千分之四（详见 \file{src/utils/hit-fonts.dtx} 里字号那张对照表）。
% 下文凡是数值一律带 \texttt{bp} 或 \texttt{pt} 写清楚，「磅」只在引 Word 界面原话时出现，
% 那时候指的是 \texttt{bp}。
%
% \emph{一、版心只由页边距定。} 宽是纸宽减左右边距，高是纸高减上下边距。
% 文档网格不参与，页眉页脚也不占版心。
%
% \emph{二、行距与每页行数互推，行距为准。} Word 存的是行距，单位是缇
% （\texttt{twip}，微软文档里的中译，一缇是 $1/20$\,\texttt{bp}）；每页行数是显示出来的
% 推导值。给行数时 Word 拿版心高的缇数除以行数，往下截而不是四舍五入：版心高
% 229\,mm 是 12982.68 缇，每页 30 行给出 432 缇也就是 21.6\,bp，量出来的十个行距一个不差
% 都是 90 个设备单位。反过来，范例存的是 391 缇也就是 19.55\,bp，Word 显示的
% 「每页 33 行」是 12982.68 除以 391 往下截得来的。两个都给且对不上时以行距为准。
%
% \emph{三、多倍行距是网格行距的倍数。} 与字体的单倍行距无关。1.25、1.5、2 倍三份对照
% 文档都对得上。
%
% \emph{四、版心顶到首行基线的距离等于首行基线在自己行盒里的落点。} 落点怎么算见
% \file{src/utils/hit-format.dtx} 的「换算」一节，对齐网格与否走两条式子。本模块
% 不管这件事：它只知道页面设置，不知道正文段落是什么字号、多少行距。
%
% Word 把每条基线四舍五入到三百分之一英寸（0.24\,bp）再落笔，累加本身是精确的。所以
% 范例里相邻行距是 19.44 与 19.68 交替，平均下来才是 19.55。这层取整是 Word 内部的，
% 用 Word 自带的「导出为 PDF」也一样，我们这边不复现。
%
% \emph{五、字距。} \cs{CJKglue} 是字距减字号。范例的 12.45\,bp 减小四 12\,bp 得
% 0.45\,bp。每行字数与字距的互推关系同第二条。
%
% \emph{六、页眉页脚。} 页眉那个距离量到页眉块的顶，页脚那个量到页脚块的底。范例里
% 页眉 3\,cm 而页眉基线在 9.04\,bp 之下，页脚 2.3\,cm 而页脚基线在 2.00\,bp 之上，两处
% 都是小五。页眉块的底超过上边距时正文整体下移同样的量，这条在「页眉 4\,cm」那份
% 对照文档上验过。
%
% \subsection{参数与产物}
%
% 尺寸类的键都收长度：\option{margin-top}、\option{margin-bottom}、\option{margin-left}、
% \option{margin-right}、\option{header}、\option{footer}、\option{line-pitch}、\option{char-pitch}、
% \option{size}。\option{lines-per-page} 与 \option{chars-per-line} 收整数。纸张默认 A4，
% 用 \option{paper-width}、\option{paper-height} 改。
%
% 另有两个键是把类自己的页眉页脚接到 Word 的位置上的：\option{head} 是 \cs{headheight}，
% \option{foot-baseline} 是页脚基线高出页脚线的量。\option{head} 不给就取调用当时的
% \cs{headheight}，不给等于零，那样会把类自己的页眉高度抹掉；\option{foot-baseline} 的
% 默认值是零，也就是页脚基线正好落在页脚线上。
%
% \emph{七、一行只排整数个字，余量留在右边。} 字符网格是一排死格子，一行装得下
% 几个字就是几个，装不满的余量原样留在右边不摊开。范例版心宽 425.18\,bp、字距
% 12.4544\,bp，装 34 个字占 423.45\,bp，右边留 1.73\,bp 空着。两端对齐是对齐到
% 这 34 个格子的右边界，不是对齐到版心右边界——范例里首行缩进两格的那一行只有
% 31 个字，被拉开填满剩下的 32 格，末字落点与满行的 34 字行分毫不差。
%
% 算完当场调用 \cs{geometry}，另外把四个量放进全局长度供调用方取用：
% \cs{g\_hit\_msword\_lead\_dim}（网格行距）、\cs{g\_hit\_msword\_char\_dim}（网格字距）、
% \cs{g\_hit\_msword\_glue\_dim}（\cs{CJKglue}，字距减字号）、
% \cs{g\_hit\_msword\_ecglue\_dim}（\cs{CJKecglue}，汉字与西文之间那一段，
% 一个增量加 0.225 格；Word 那边两侧合计 0.45 格，四个字号量下来除以格是常数、
% 除以字号在漂，所以按格算）、
% \cs{g\_hit\_msword\_slack\_dim}（第七条那点余量，交给 \cs{rightskip}）。
% 前两个是 \file{src/utils/hit-format.dtx} 里「一行」「一字」这两个单位的取值来源。
%
% ^^A ======================================================================
% ^^A Module shared-msword: Word 页面设置换算（各个类共用）
% ^^A ======================================================================
%    \begin{macrocode}
%<*shared-msword>
%    \end{macrocode}
%
% 缇是 Word 的内部长度单位，第二条规则的取整在这个单位上做。一缇是 $1/20$\,\texttt{bp}，
% 写成 \texttt{0.05pt} 会让每页 30 行算出 434 缇而不是 432 缇。
%
% \texttt{pt} 换 \texttt{bp} 是乘 $7200/7227$，再乘 20 得缇，两步并成一个有理数
% \texttt{144000/7227}。不把 \texttt{0.05bp} 存成长度再乘份数：长度只能到
% \texttt{sp}（$1/65536$\,pt），一个缇存进去差 0.09\,sp，乘上四百多倍就差到
% $0.001$\,\texttt{bp}，排出来的行距对不上整数缇。
%    \begin{macrocode}
\int_const:Nn \c__hit_msword_twip_int { 144000 }
\int_const:Nn \c__hit_msword_point_int { 7227 }
%    \end{macrocode}
%
% 输出的三个量。
%    \begin{macrocode}
\dim_new:N \g_hit_msword_lead_dim
\dim_new:N \g_hit_msword_char_dim
\int_new:N \g_hit_msword_cells_int
%    \end{macrocode}
%
%    \begin{macrocode}
\dim_new:N \l__hit_msword_paper_width_dim
\dim_new:N \l__hit_msword_paper_height_dim
\dim_new:N \l__hit_msword_top_dim
\dim_new:N \l__hit_msword_bottom_dim
\dim_new:N \l__hit_msword_left_dim
\dim_new:N \l__hit_msword_right_dim
\dim_new:N \l__hit_msword_header_dim
\dim_new:N \l__hit_msword_footer_dim
\dim_new:N \l__hit_msword_header_height_dim
\dim_new:N \l__hit_msword_rule_width_dim
\dim_new:N \l__hit_msword_rule_text_dim
\dim_new:N \l__hit_msword_foot_shift_dim
\dim_new:N \l__hit_msword_foot_shift_x_dim
\dim_new:N \g_hit_msword_foot_shift_x_dim
\dim_new:N \l__hit_msword_lead_dim
\dim_new:N \l__hit_msword_char_dim
\dim_new:N \l__hit_msword_text_width_dim
\dim_new:N \l__hit_msword_text_height_dim
\int_new:N \l__hit_msword_lines_int
\int_new:N \l__hit_msword_chars_int
\int_new:N \l__hit_msword_cells_int
\tl_new:N \g__hit_msword_geometry_tl
%    \end{macrocode}
%
% \option{line-pitch} 与 \option{lines-per-page} 谁都可以不给，所以两个的初值是 0，
% 拿来当「没给」的哨兵。\option{char-pitch} 与 \option{chars-per-line} 同理。
%    \begin{macrocode}
\keys_define:nn { hit / msword }
  {
    paper-width    .dim_set:N = \l__hit_msword_paper_width_dim ,
    paper-width    .initial:n = { 210mm } ,
    paper-height   .dim_set:N = \l__hit_msword_paper_height_dim ,
    paper-height   .initial:n = { 297mm } ,
    margin-top     .dim_set:N = \l__hit_msword_top_dim ,
    margin-bottom  .dim_set:N = \l__hit_msword_bottom_dim ,
    margin-left    .dim_set:N = \l__hit_msword_left_dim ,
    margin-right   .dim_set:N = \l__hit_msword_right_dim ,
    header         .dim_set:N = \l__hit_msword_header_dim ,
    footer         .dim_set:N = \l__hit_msword_footer_dim ,
    header-rule-width     .dim_set:N = \l__hit_msword_rule_width_dim ,
    header-rule-width     .initial:n = { 2.25bp } ,
    header-rule-from-text .dim_set:N = \l__hit_msword_rule_text_dim ,
    header-rule-from-text .initial:n = { 1bp } ,
    footer-shift-vertical .dim_set:N = \l__hit_msword_foot_shift_dim ,
    footer-shift-vertical .initial:n = { 0pt } ,
    line-pitch     .dim_set:N = \l__hit_msword_lead_dim ,
    line-pitch     .initial:n = { 0pt } ,
    char-pitch     .dim_set:N = \l__hit_msword_char_dim ,
    char-pitch     .initial:n = { 0pt } ,
    lines-per-page .int_set:N = \l__hit_msword_lines_int ,
    lines-per-page .initial:n = { 0 } ,
    chars-per-line .int_set:N = \l__hit_msword_chars_int ,
    chars-per-line .initial:n = { 0 } ,
    footer-shift-horizontal .dim_set:N = \l__hit_msword_foot_shift_x_dim ,
    footer-shift-horizontal .initial:n = { 0pt } ,
  }
%    \end{macrocode}
%
% \cs{\_\_hit\_msword\_pitch:NNn}\meta{目标}\meta{版心尺寸}\meta{份数} 是第二条规则：
% 目标已经有值就不动，否则按份数除，在缇上往下截。
%    \begin{macrocode}
\cs_new_protected:Npn \__hit_msword_pitch:NNn #1#2#3
  {
    \bool_lazy_and:nnT
      { \dim_compare_p:nNn {#1} = { 0pt } }
      { \int_compare_p:nNn {#3} > { 0 } }
      {
        \dim_set:Nn #1
          {
            \fp_to_decimal:n
              {
                floor
                  (
                    \dim_to_decimal:n {#2}
                    * \c__hit_msword_twip_int / \c__hit_msword_point_int / #3
                  )
                / 20
              }
            bp
          }
      }
  }
%    \end{macrocode}
%
% 主命令。键是局部的，整段套在组里，出组时自动复位，下一次调用不受这一次影响。三个
% 产物与传给 \pkg{geometry} 的键值表跨组，所以是全局赋值。
%    \begin{macrocode}
\cs_new_protected:Npn \hit@msword@layout:n #1
  {
    \group_begin:
      \keys_set:nn { hit / msword } {#1}
%    \end{macrocode}
%
% \changes{v3.2a}{2026/08/16}{页眉盒高由页眉行高、\option{header}/\option{rule}/%
%   \option{from-text} 与线宽算出来，不再手写}
% 页眉盒高不是独立的量，它是三段之和：页眉那一行的自然行高、文字到横线的距离、
% 横线本身。Word 那边横线是页眉段落的下边框，位置就是这么定的——跟着文字走，
% 隔一个 \texttt{w:space}，再画线。所以这里照着算，别人手写一个数就失去了这层关系。
%
% 横线是 \texttt{thinThickSmallGap}：粗线取 \option{width}（Word 的 \texttt{w:sz}），
% 间隙与细线各 0.75\,磅，比例由线型定义写死，Word 也不让改。
%
% 页眉一律小五，所以行高直接取小五。哪天页眉字号可配了，这一句跟着换。
%    \begin{macrocode}
      \dim_set:Nn \l__hit_msword_header_height_dim
        {
          \c__hit_format_natural_tl \c__hit_font_size_xiaowu_dim
          + \l__hit_msword_rule_text_dim
          + \l__hit_msword_rule_width_dim + 1.5bp
        }
      \dim_gset_eq:NN \g_hit_msword_foot_shift_x_dim
        \l__hit_msword_foot_shift_x_dim
      \dim_set:Nn \l__hit_msword_text_width_dim
        {
          \l__hit_msword_paper_width_dim
          - \l__hit_msword_left_dim - \l__hit_msword_right_dim
        }
      \dim_set:Nn \l__hit_msword_text_height_dim
        {
          \l__hit_msword_paper_height_dim
          - \l__hit_msword_top_dim - \l__hit_msword_bottom_dim
        }
      \__hit_msword_pitch:NNn \l__hit_msword_lead_dim
        \l__hit_msword_text_height_dim { \l__hit_msword_lines_int }
      \__hit_msword_pitch:NNn \l__hit_msword_char_dim
        \l__hit_msword_text_width_dim { \l__hit_msword_chars_int }
      \dim_compare:nNnT { \l__hit_msword_lead_dim } = { 0pt }
        { \msg_error:nn { hithesis } { msword-no-pitch } }
      \dim_gset_eq:NN \g_hit_msword_lead_dim \l__hit_msword_lead_dim
      \dim_gset_eq:NN \g_hit_msword_char_dim \l__hit_msword_char_dim
%    \end{macrocode}
%
% \changes{v3.2a}{2026/08/16}{余量的算法挪进 \file{src/utils/hit-format.dtx}，
%   这里只出「一行放得下几个格」}
% 一行放得下几个格，只跟版心宽与字距有关，是版面的量。余量还要减掉正文字号，
% 那个数由 \option{format}/\option{body}/\option{size} 定，所以算法挪去 format，
% 见 \cs{hit@format@slack:}。
% 字距没给就没有网格，格数是零。
%    \begin{macrocode}
      \int_zero:N \l__hit_msword_cells_int
      \dim_compare:nNnF { \l__hit_msword_char_dim } = { 0pt }
        {
          \int_set:Nn \l__hit_msword_cells_int
            {
              \fp_to_int:n
                {
                  floor
                    (
                      \dim_to_decimal:n { \l__hit_msword_text_width_dim }
                      / \dim_to_decimal:n { \l__hit_msword_char_dim }
                    )
                }
            }
        }
      \int_gset_eq:NN \g_hit_msword_cells_int \l__hit_msword_cells_int
      \tl_gset:Nx \g__hit_msword_geometry_tl
        {
          paperwidth = \dim_use:N \l__hit_msword_paper_width_dim ,
          paperheight = \dim_use:N \l__hit_msword_paper_height_dim ,
          left = \dim_use:N \l__hit_msword_left_dim ,
          right = \dim_use:N \l__hit_msword_right_dim ,
          top = \dim_use:N \l__hit_msword_top_dim ,
          bottom = \dim_use:N \l__hit_msword_bottom_dim ,
          head = \dim_use:N \l__hit_msword_header_height_dim ,
          headsep = \dim_eval:n
            {
              \l__hit_msword_top_dim
              - \l__hit_msword_header_dim - \l__hit_msword_header_height_dim
            } ,
          footskip = \dim_eval:n
            {
              \l__hit_msword_bottom_dim
              - \l__hit_msword_footer_dim + \l__hit_msword_foot_shift_dim
            } ,
        }
    \group_end:
    \exp_args:NV \geometry \g__hit_msword_geometry_tl
  }
%    \end{macrocode}
%
%    \begin{macrocode}
%</shared-msword>
%    \end{macrocode}
%
% \Finale
